\chapter{Conclusion}
\label{chap:conclusion}

Ce stage avait pour objectif principal de spécialiser LANCE par une architecture
HMoE agentique, avec une empreinte d'entraînement compatible avec des ressources
limitées. La solution partage un modèle de base quantifié entre plusieurs
adaptateurs QLoRA associés aux rôles de la campagne. Un routeur déterministe
sélectionne l'expert attendu selon la phase et le serveur expose l'ensemble
derrière une interface commune. La mission a couvert la préparation des traces,
l'entraînement, le chargement des adaptateurs et leur intégration au cycle
d'action de l'agent.

Le pipeline et le harness ont été consolidés pour rendre cette spécialisation
évaluable. Le pipeline encadre les transitions, restreint les outils et valide
les livrables. Le harness déploie des scénarios reproductibles, associe les
résultats à une vérité terrain et mesure la qualité des campagnes. Les travaux
sur l'API, le suivi des exécutions, l'automatisation de l'infrastructure et les
métriques participent donc à un même protocole expérimental. Le graphe pondéré
et le cône d'impact apportent en complément un signal de priorité lorsque l'état
du réseau évolue.

Les résultats provisoires montrent que l'agent peut achever des campagnes sur
plusieurs scénarios, mais aussi que la production d'un constat ne garantit pas
un chemin d'attaque entièrement prouvé. Cet écart confirme l'intérêt d'évaluer
le modèle au niveau de son exécution plutôt qu'à partir de la seule
vraisemblance de ses réponses.

Plusieurs limites demeurent. Les adaptateurs dépendent de la diversité et de la
correction des traces. Le routage par phase est interprétable, mais ne combine
pas plusieurs spécialités aux frontières des rôles. Les poids du graphe restent
heuristiques et les premiers résultats ne constituent pas encore une comparaison
définitive avec un modèle généraliste.

La suite devra exécuter les configurations sur un jeu de scénarios identique,
répéter les campagnes et analyser leur qualité, leurs preuves et leurs coûts.
Une boucle d'amélioration pourra ensuite soumettre les erreurs à une revue
humaine avant d'enrichir les données. Ce projet m'a finalement appris qu'un
modèle spécialisé ne peut être jugé indépendamment de l'architecture qui
encadre ses actions et du protocole qui rend ses résultats vérifiables.
